% Appendix D --- hypergeometric Laplace transforms of the kernels.  The
% V_inf^(k) derivation and the PGF coefficient extraction live in the
% distribution-theory chapter's technical appendix.

\section{Technical derivations}
\label{p:app:technical}

\subsection{Hypergeometric Laplace transforms of the kernels}
\label{p:app:transforms}

Write $s:=r/\theta$ and recall $\theta=\lambda(a-b)$. Both transforms
reduce, under $v=\ee^{-\theta t}$ (so $\dd t=-\dd v/(\theta v)$, $v:1\to0$),
to moments of the form $\int_0^1 v^{s-1}(1-v)^n(A+Bv)^{-m}\dd v$, which are
hypergeometric. For the survival kernel, from \cref{p:eq:IDIhat},
\[
\Ifix(t)=\frac{(a-b)^2w}{(B+Aw)(aw-b)},\qquad w=\ee^{\theta t},
\]
and in $v=1/w$,
\[
\Lap{\Ifix}(r)=\frac{a-b}{\theta}\int_0^1 v^s
\Bigl[\frac{B}{A+Bv}+\frac{b}{a-bv}\Bigr]\dd v,
\]
using the partial-fraction decomposition
$(a-b)/[(A+Bv)(a-bv)]=B/(A+Bv)+b/(a-bv)$ --- both sides multiplied out give
$(a-b)^2$ in the numerator after using $Ab+aB=a-b$. Each term is a standard
hypergeometric integral,
$\int_0^1 v^s/(1-zv)\dd v=(s+1)^{-1}\Fhyp(1,s+1;s+2;z)$, giving
\begin{equation}
\Lap{\Ifix}(r)=\frac{1}{\lambda(s+1)}
\Biggl[
\Psi(A,B,s)
+\frac{b}{a}\,\Fhyp\Bigl(1,s+1;s+2;\frac{b}{a}\Bigr)
\Biggr],
\label{p:eq:LIhat}
\end{equation}
where, so that the hypergeometric series always converges,
\begin{equation}
\Psi(A,B,s)=
\begin{cases}
\Fhyp\bigl(1,s+1;s+2;-\frac{A}{B}\bigr), & B\ge A,\\[0.4em]
\frac{B}{A}\,\Fhyp\bigl(1,s+1;s+2;-\frac{B}{A}\bigr), & A>B,
\end{cases}
\label{p:eq:Psi}
\end{equation}
the two branches being equal by the identity
$(B/A)\Fhyp(1,s+1;s+2;-B/A)=\Fhyp(1,s+1;s+2;-A/B)$ --- both sides equal
$(s+1)B\int_0^1 v^s/(A+Bv)\dd v$. (The single-argument form
$(B/A)\Fhyp(1,s+1;s+2;-B/A)$, valid by analytic continuation for all
$A,B>0$, has a convergent defining series only when $B\le A$, that is
$\delta\ge\lambda-\mu$.) For the release kernel, $g=\delta K$ with
$K=[1+2\lambda(1-I)/\delta]J$; the same substitution and termwise
integration give
\begin{equation}
\delta\Lap{K}(r)=\frac{\delta}{\lambda}
\Biggl[
\frac{\Fhyp(1,s;s+2;z)}{A(s+1)}
+\frac{2\,\Fhyp(2,s;s+3;z)}{A^2(s+1)(s+2)}
\Biggr],
\qquad z=-\frac{B}{A},
\label{p:eq:dLK}
\end{equation}
again with the understanding that for $B>A$ the defining series of the
right-hand side must be evaluated through its analytic continuation, or
through the integral representation from which it came. Both formulae reduce
correctly at $s=0$: with $\Fhyp(1,1;2;z)=-\log(1-z)/z$, \cref{p:eq:LIhat}
gives $\Lap{\Ifix}(0)=\lambda^{-1}\log(a/(a-1))=\mean{T_{\rm prod}}$, and
\cref{p:eq:dLK} gives $\delta\Lap{K}(0)=(\delta/\lambda)(1/A+1/A^2)
=aB/A=V_\infty$, using $AB=\delta/\lambda$. Both reductions were verified
numerically against direct quadrature, as was \cref{p:eq:dLK} for general
$r$, to relative error $\sim10^{-8}$ in the convergent regime. Where the
transforms are needed as functions of time rather than of $r$, the standard
numerical inversions apply
\cite{stehfest1970algorithm,gaver1965observations}.
% NEEDS-REF: a standard Volterra-quadrature reference for the direct
% solution of the renewal system

The companion derivations --- the lifetime yield $V_\infty^{(k)}$ and the
generating-function coefficient extraction used by the chained model --- are
in \ChDist{}'s technical appendix.

\FloatBarrier
